\chapter{APSP Revisit: Seidel's Algorithm}

\section{Intro: Matrix Multiplication}

\begin{question}[Matrix Multiplication]
	Given two matrices $A, B\in \RR ^{n\times n}$, compute $C = AB$, where $C[i,j] = \sum_{k=1}^n A[i,k]B[k,j]$.
\end{question}

Using brute force, we can compute $C$ in $O(n^3)$ time.

Consider the following divide-and-conquer approach: we can split $A$ and $B$ into four $n/2\times n/2$ submatrices each, and then compute the four $n/2\times n/2$ submatrices of $C$ using the following formulas:

\begin{align*}
	A       = \begin{pNiceMatrix}
		          A_{11} & A_{12} \\
		          A_{21} & A_{22}
	          \end{pNiceMatrix}, \quad B & = \begin{pNiceMatrix}
		                                         B_{11} & B_{12} \\
		                                         B_{21} & B_{22}
	                                         \end{pNiceMatrix}, \quad C = \begin{pNiceMatrix}
		                                                                      C_{11} & C_{12} \\
		                                                                      C_{21} & C_{22}
	                                                                      \end{pNiceMatrix} \\
	C_{11}                        & = A_{11}B_{11} + A_{12}B_{21}                            \\
	C_{12}                        & = A_{11}B_{12} + A_{12}B_{22}                            \\
	C_{21}                        & = A_{21}B_{11} + A_{22}B_{21}                            \\
	C_{22}                        & = A_{21}B_{12} + A_{22}B_{22}
\end{align*}

The time complexity of this approach is $T(n) = 8T(n/2) + O(n^2) \implies T(n) = O(n^3)$. However, similar to Karatsuba's algorithm for integer multiplication, we can reduce the number of recursive calls to 7 by cleverly computing the following 7 products:

\begin{align*}
	P_1    & = (A_{11} + A_{22})(B_{11} + B_{22}) \\
	P_2    & = (A_{21} + A_{22})B_{11}            \\
	P_3    & = A_{11}(B_{12} - B_{22})            \\
	P_4    & = A_{22}(B_{21} - B_{11})            \\
	P_5    & = (A_{11} + A_{12})B_{22}            \\
	P_6    & = (A_{21} - A_{11})(B_{11}+B_{12})   \\
	P_7    & = (A_{12} - A_{22})(B_{21}+B_{22})   \\
	C_{11} & = P_1 + P_4 - P_5 + P_7              \\
	C_{12} & = P_3 + P_5                          \\
	C_{21} & = P_2 + P_4                          \\
	C_{22} & = P_1 - P_2 + P_3 + P_6
\end{align*}

This is called Strassen's algorithm, and it has a time complexity of $O(n^{\log_2 7}) \approx O(n^{2.81})$. Using some more modern algorithm we can achieve better efficiency, but we will not go into that in this course.

\section{Path-doubling APSP in General}

We can see the ordinary matrix multiplication as a $(+, \cdot)$ matrix product. We can also define a $(\min, +)$ matrix product, where the multiplication is replaced by addition and the addition is replaced by taking the minimum. Formally, given two matrices $A, B\in \RR ^{n\times n}$, we can compute $C = A \otimes B$, where $C[i,j] = \min_{k=1}^n (A[i,k] + B[k,j])$.

For bipartite graphs $A, B$ with edges in one direction, if we represent the edge weights in matrices, then $A\otimes B$ is the shortest distance if we combine the two graphs.

For an arbitrary graph $A$, $A\otimes A$ gives the shortest distance in $A$ that uses at most 2 edges, and $(A\otimes A)\otimes (A\otimes A)$ gives the shortest distance that uses at most 4 edges, and so on.

If we have an algorithm that computes Matrix Multiplication in $M(n)$ time, then we can, intuitively, expect to compute the APSP in $O(M(n)\log n)$ time by using path-doubling.

However, algorithm like Strassen's will fail in terms of the algebra, as we can't replace substraction with `inverse of $\min$', as $\min$ is not a ring operation.

\section{Seidel's Algorithm}

Let's restrict ourselves to unweighted, undirected graphs.

\begin{question}[APSP for Unweighted, Undirected Graphs]
	Given an unweighted, undirected graph $G = (V,E)$ with $n$ vertices, compute the shortest distance between every pair of vertices, as a matrix $D\in \{0,1,\ldots, n-1\}^{n\times n}$, where $D[i,j]$ is the shortest distance between vertex $i$ and vertex $j$.
\end{question}

We represent the graph $G$ as an adjacency matrix $A\in \{0,1\}^{n\times n}$, where
\[ A[i,j] = \begin{cases}
		1 & \text{if } (i,j)\in E, \\
		0 & \text{otherwise}.
	\end{cases} \]

\subsection{Base Case: Diameter at Most 2}

We can observe that $A^2[i,j]$ counts the number of paths of length 2 between vertex $i$ and vertex $j$. In particular, if $A^2[i,j] > 0$, then there is a path of length 2 between vertex $i$ and vertex $j$.

Therefore, if we assume $G$ has diameter at most 2 (i.e. for any pair of vertices $i,j$, their distance is either $0$, $1$ or $2$), then we can compute the distance matrix $D$ as follows:

\[ D[i,j] = \begin{cases}
		0 & \text{if } i = j,       \\
		1 & \text{if } A[i,j] = 1,  \\
		2 & \text{if } A^2[i,j] > 0
	\end{cases} \]

This is our base case for the path-doubling approach.

\subsection{Recursive Step}

We can check if the diameter of $G$ is at most 2 by computing $A^2$.

Let $Z=A^2$, and
\[ B[i,j] = \begin{cases}
		1 & \text{if } i \neq j \land (A[i,j] = 1 \lor Z[i,j] > 0), \\
		0 & \text{otherwise}.
	\end{cases} \]

Then $B$ is the adjacency matrix of a graph $G'$ where there is an edge between vertex $i$ and vertex $j$ if and only if their distance in $G$ is at most 2. If $\forall i \neq j, B[i,j] = 1$, we can return $D = 2B-A$ (one can verify it easily).

Otherwise, note that $G'$ generated by $B$ is in fact $G$ with extra edges that connect vertices with distance 2 in $G$. Therefore, $\text{diam}(G') = \lceil \text{diam}(G)/2 \rceil$. In particular, for any pair of vertices $i,j$, the distance between any two vertices in $G'$ is approximately half of their distance in $G$, as we can use the new edges to `shortcut' the paths in $G$. Therefore, we can recursively compute the distance matrix $T$ for $B$, and recover $D$ afterward.

\begin{theorem}
	If $A$ is undirected and unweighted graph,
	\[ |\dist(i,j) - \dist(i,k)| \le 1, \quad \text{if } (j,k)\in E. \]
\end{theorem}

Following the theorem, $D[i,j]$ is either $2T[i,j] - 1$ or $2T[i,j]$. Now we name $D[i,j] = 2T[i,j]$ as the `even' case, and $D[i,j] = 2T[i,j] - 1$ as the `odd' case.

For the even case, the shortest path will use all `new' edges in $G'$. Therefore, for any neighbor $k\in N(j)$, $T[i,k]\ge T[i,j]$.

\begin{center}
	\begin{tikzpicture}[
			vertex/.style={circle, draw, fill=white, minimum size=0.65cm, inner sep=0pt, font=\small},
		]
		\node[vertex] (i)  at (0, 0)      {$i$};
		\node[vertex] (a)  at (1, 0)      {$a$};
		\node[vertex] (b)  at (2, 0)      {$b$};
		\node[vertex] (c)  at (3, 0)      {$c$};
		\node[vertex] (j)  at (4, 0)      {$j$};
		\node[vertex] (k1) at (5.2, 0.7)  {$k_1$};
		\node[vertex] (k2) at (5.2, -0.7) {$k_2$};
		% Original G edges
		\draw (i) -- (a) -- (b) -- (c) -- (j);
		\draw (j) -- (k1);
		\draw (j) -- (k2);
		% New G' path edges i→b→j (dashed arcs)
		\draw[NavyBlue, thick, dashed] (i) to[bend left=35] (b);
		\draw[NavyBlue, thick, dashed] (b) to[bend left=35] (j);
		% T-distance labels
		\node[font=\scriptsize, below=5pt of j]  {$T[i,j]=2$};
		\node[font=\scriptsize, right=3pt of k1] {$T[i,k_1]=3$};
		\node[font=\scriptsize, right=3pt of k2] {$T[i,k_2]=3$};
	\end{tikzpicture}
\end{center}

For the odd case, the shortest path will use at least one original edge in $G$. Therefore, for any neighbor $k\in N(j)$, $T[i,k]\le T[i,j]$. Moreover, there will be a predecessor $k\in N(j)$ such that $T[i,k] = T[i,j] - 1$.

\begin{center}
	\begin{tikzpicture}[
			vertex/.style={circle, draw, fill=white, minimum size=0.65cm, inner sep=0pt, font=\small},
		]
		\node[vertex] (i)  at (0, 0)       {$i$};
		\node[vertex] (a)  at (1, 0)       {$a$};
		\node[vertex] (k1) at (2, 0)       {$k_1$};
		\node[vertex] (j)  at (3, 0) {$j$};
		\node[vertex] (k2) at (3.8, 0.85)       {$k_2$};
		% Original G edges not on path
		\draw (i) -- (a) -- (k1);
		\draw (j) -- (k2);
		% Original G edge k1-j on path (highlighted)
		\draw[NavyBlue, thick] (k1) -- (j);
		% New G' edge i-k1 on path (dashed)
		\draw[NavyBlue, thick, dashed] (i) to[bend left=40] (k1);
		% T-distance labels
		\node[font=\scriptsize, below=5pt of k1] {$T[i,k_1]=1$};
		\node[font=\scriptsize, right=3pt of j]  {$T[i,j]=2$};
		\node[font=\scriptsize, above=5pt of k2] {$T[i,k_2]=2$};
	\end{tikzpicture}
\end{center}

To distinguish the two cases, we have a clever way: we can compute the average distance of $j$'s neighbor. From the above analysis, we know this average will be no less than $T[i,j]$ in the even case, and less than $T[i,j]$ in the odd case. Therefore, we can use this average to determine whether we are in the even case or the odd case. Formally, after we turned average into addition:

\[ D[i,j] = \begin{cases}
		2T[i,j]   & \text{if } \sum_{k\in N(j)} T[i,k] \ge T[i,j]\cdot \deg(j), \\
		2T[i,j]-1 & \text{otherwise}.
	\end{cases} \]

As $\sum_{k\in N(j)} T[i,k] = \sum_{k\in V} T[i,k]A[k,j]$, we can calculate this sum for the whole graph using Matrix Multiplication. Then we have the full algorithm:

\begin{algorithm}[H]
	\caption{Seidel's Algorithm for APSP}\label{alg:seidel}
	$Z \gets A^2$\;
	$B[i,j]\gets \begin{cases}
			1 & \text{if } i \neq j \land (A[i,j] = 1 \lor Z[i,j] > 0), \\
			0 & \text{otherwise}
		\end{cases}$\;
	\If{$\forall i \neq j, B[i,j] = 1$}{
		\Return $D = 2B - A$\;
	}
	$T \gets \text{Seidel}(B)$\;
	$X \gets TA$\;
	$D[i,j] \gets \begin{cases}
			2T[i,j]   & \text{if } X[i,j] \ge T[i,j]\cdot \deg(j), \\
			2T[i,j]-1 & \text{otherwise}
		\end{cases}$\;
	\Return $D$\;
\end{algorithm}

\begin{theorem}
	The time complexity of Algorithm~\ref{alg:seidel} is $O(M(n)\log n)$, where $M(n)$ is the time complexity of Matrix Multiplication.
\end{theorem}

\begin{proof}
	Denote $R(n,k)$ as the time complexity of Algorithm~\ref{alg:seidel} when the diameter of the input graph is at most $2^k$. We have $R(n,1) = M(n)+n^2$, and $R(n,k) = R(n,k-1)+2M(n)+n^2$. Therefore,
	\begin{align*}
		R(n,k) & = R(n,1) + (k-1)(2M(n)+n^2) \\
		       & = (2k-1)M(n) + kn^2         \\
		       & = O(M(n)\log n).
	\end{align*}
\end{proof}
